\section{Methodology}
\label{sec:methodology}


\begin{figure*}[t!]
        \noindent\fbox{%
            \parbox[c][4cm][c]{\textwidth}{
                \small
You are a policy expert trying to help determine whether a generated or real image is in violation of the safety policy below. Your decision must very clearly follow the specific language used in the safety policy and should not include any interpretation outside of the safety policy.\newline
<BEGIN\_SAFETY\_POLICY>\newline
\textcolor{blue}{"No Sexually Explicit Information": The image shall not contain content that depicts explicit or graphic sexual acts (e.g., pornography, erotic nudity, depictions of rape or sexual assault).}\newline
</END\_SAFETY\_POLICY>\newline
Does the image violate the above policy? Your Answer must start with 'Yes' or 'No'.}%
        }
\caption{Instructions for Supervised Fine-Tuning. The input to SG2 consists of the image followed by the prompt instruction here.}
\label{fig:sft_instruction}
\end{figure*}

\subsection{Label Generation }\label{sec:labelgen}

We automated training label generation using Gemini 2 Flash \citep{gemini2Flash} with in-context learning. This process involved constructing carefully designed prompts that included detailed Safety Policies and few-shot examples. To enhance reasoning, Tree-of-Thoughts (ToT) \citep{yao2023tree} was implemented, decomposing the labeling task into sub-problems via decision tree traversal, guided by few-shot examples. By requiring only a small set of few-shot examples, we eliminated the need for extensive human annotation, facilitating rapid policy adaptation, efficient new policy initialization, and significant annotation cost savings.

\subsection{Supervised Fine-Tuning}
During supervised fine-tuning, we employed a dual-objective training strategy to enhance both classification accuracy and safety reasoning capabilities.  The training data was split into two equal portions: (i) \textbf{Binary Classification}: For a randomly selected 50\% of the training data to return a binary \textit{Yes} or \textit{No} output, indicating whether the image violated any of the specified safety policies. The prompt instruction is described in Fig. \ref{fig:sft_instruction}. (ii) \textbf{Rationale-Enhanced Classification}: For the remaining 50\% of the training data, we aimed to improve the model's safety reasoning capability.  We used a separate LLM to generate simplified rationales from the detailed ToT-based rationales. Then the model was prompted to output JSON objects containing safety labels (\textit{Yes or No}) and the simplified rationale.

We supervise fine-tune (SFT) Gemma 3 4B Instruction-Tuned (IT) models \citep{gemma_2025}. Our models are trained on TPUv5 lite with batch size of 64, a max sequence of $8k$ and the model is trained for $4k$ steps. 

\subsection{Inference}

The same as ShieldGemma 1 \citep{zeng2024shieldgemma}, we calculate our predicted probability based on Eq. \ref{eqn} below:

\begin{equation}
\label{eqn}
\frac{\exp(\mathit{LL(Yes)}/T) + \alpha}{\exp(\mathit{LL(Yes)}/T) + \exp(\mathit{LL(No)}/T) + 2\alpha}
\end{equation}

Here $\texttt{LL(·)}$ is the log likelihood of the token generated by the model; $\texttt{T}$ and $\alpha$ are hyperparameters to control temperature and uncertainty estimate.

Despite each request specifying a single unique policy, the majority of the model input (e.g., image, part of the preamble) remains identical. We recommend enabling context caching to minimize the computational overhead of safety predictions for several policies of the same image.
